\hypertarget{appendix-b9afba7f63e5}{}\label{referencestart:appendix-b9afba7f63e5}
\section{Stable quadratic root evaluation}\label{appendix-b9afba7f63e5-stable-quadratic-root-evaluation}

\hypertarget{appendix-b9afba7f63e5-stable-quadratic-root-evaluation}{}

\textbf{Curriculum depth:} Developed · \textbf{First seen:} Phase 10 · \textbf{Latest development:} Phase 10

\begin{allmightroledefinition}

\subsection{Definition and mechanism}\label{appendix-b9afba7f63e5-definition-and-mechanism}

\hypertarget{appendix-b9afba7f63e5-definition-and-mechanism}{}

The direct quadratic formula can lose nearly all significant digits in one root when it subtracts nearly equal values. For real coefficients with nonzero leading coefficient and a valid nonnegative discriminant, compute a signed large-magnitude numerator first: \(q=-\frac12(b+\operatorname{copysign}(\sqrt{b^2-4ac},b))\), then obtain roots as \(q/a\) and \(c/q\). The second root follows from the product relation \(x_1x_2=c/a\). This avoids cancellation in the small root, but it does not prevent overflow in the discriminant or cover degenerate coefficients, nonfinite inputs, or \(q=0\); a robust solver also scales inputs and handles those cases explicitly. For example, the small root of \(x^2+10^8x+1=0\) is near \(-10^{-8}\), exactly the root most endangered by direct subtraction.

\end{allmightroledefinition}

\subsection{Prerequisites}\label{appendix-b9afba7f63e5-prerequisites}

\hypertarget{appendix-b9afba7f63e5-prerequisites}{}

\begin{itemize}
\tightlist
\item
  \href{../01-Bits-Programs-and-Reliable-Systems-P01/main.pdf\#appendix-d7f8ba961670}{Stable reformulation (Bits, Programs, and Reliable Systems, Part 1, Reference notes: Stable reformulation)}
\item
  \href{../01-Bits-Programs-and-Reliable-Systems-P01/main.pdf\#appendix-bc44f7b42f42}{Unit roundoff (Bits, Programs, and Reliable Systems, Part 1, Reference notes: Unit roundoff)}
\item
  \href{../01-Bits-Programs-and-Reliable-Systems-P01/main.pdf\#appendix-8f6a01f284a7}{Conditioning and singularity (Bits, Programs, and Reliable Systems, Part 1, Reference notes: Conditioning and singularity)}
\item
  \hyperref[appendix-aa2bf0d79639]{Algorithmic stability}
\end{itemize}

\subsection{Used by}\label{appendix-b9afba7f63e5-used-by}

\hypertarget{appendix-b9afba7f63e5-used-by}{}

\begin{itemize}
\tightlist
\item
  \hyperref[appendix-f63eb850796d]{Bracketed and open root solvers}
\end{itemize}

\subsection{Full teaching explanation}\label{appendix-b9afba7f63e5-full-teaching-explanation}

\hypertarget{appendix-b9afba7f63e5-full-teaching-explanation}{}

\begin{itemize}
\tightlist
\item
  \href{../02-Mathematics-for-Visual-Computing-P02/main.pdf\#02-foundations-ii-ch028-floating-point-answers-error-conditioning-stability-and-stable-arithmetic}{Floating-point answers: error, conditioning, stability, and stable arithmetic (Mathematics for Visual Computing, Part 2, Ch. 12)}
\end{itemize}

\subsection{Source coverage}\label{appendix-b9afba7f63e5-source-coverage}

\hypertarget{appendix-b9afba7f63e5-source-coverage}{}

Supplied lesson parts: 974.

\subsection{Future keywords}\label{appendix-b9afba7f63e5-future-keywords}

\hypertarget{appendix-b9afba7f63e5-future-keywords}{}

Vieta relation, discriminant scaling, cancellation
